\begin{ouverture}[grandeurs]{13}{Aires et angles : comparer}{Partie 3 · Grandeurs et mesures}
{Surfaces · Pavages · Centimètre carré · Écartements}
\objectif{Comparer des aires par superposition, découpage ou pavage.}
\objectif{Mesurer une aire en carreaux ou en centimètres carrés.}
\objectif{Comparer des angles et repérer les angles droits.}
\prerequis{Quadrillages et figures usuelles.}
\end{ouverture}
\pagestyle{otMain}\essentiel\fontsize{11}{15}\selectfont
\phantomsection\label{cours-13-1}
\essSection{1}{Comparer et mesurer des aires}
\begin{definition}L’aire d’une figure est la mesure de l’étendue de sa surface. Deux figures peuvent être comparées en les superposant, ou en découpant l’une pour recouvrir l’autre, sans rien mesurer.\end{definition}
\begin{savoirfaire}[Paver pour mesurer]Sur un quadrillage, on choisit un carreau comme unité d’aire : l’aire est le nombre de carreaux qui pavent la figure. Deux demi-carreaux comptent pour un carreau entier.\end{savoirfaire}
\begin{definition}Le centimètre carré, noté $\mathrm{cm}^2$, est l’aire d’un carré de 1 cm de côté. Une figure pavée par 12 de ces carrés a une aire de 12 $\mathrm{cm}^2$.\end{definition}
\begin{attention}Aire et périmètre sont deux grandeurs différentes : l’une mesure la surface, l’autre le contour. Deux figures de même périmètre peuvent avoir des aires différentes, et réciproquement.\end{attention}
\clearpage
\phantomsection\label{cours-13-2}
\essSection{2}{Comparer des angles}
\begin{definition}Un angle est l’écartement entre deux côtés qui partent d’un même sommet. L’angle de sommet B, dont les côtés passent par A et C, se note $\widehat{ABC}$ : la lettre du sommet est au milieu.\end{definition}
\label{fig-cours-13-2-cours-angles}\figureGeo{cours-angles}

\begin{savoirfaire}[Comparer avec un calque]Pour comparer deux angles, décalquer l’un et le poser sur l’autre, sommet sur sommet et un côté sur l’autre : le plus ouvert est le plus grand. La longueur dessinée des côtés ne change pas l’angle.\end{savoirfaire}
\begin{propriete}L’angle droit est l’angle du coin de l’équerre. Un angle peut être plus petit, égal ou plus grand qu’un angle droit : c’est le premier repère pour classer les angles.\end{propriete}

\clearpage\rubrique{Découverte}
\begin{blocLivre}{ex-13-1}\begin{exercice}[1]{Qui a la plus grande aire ?}
Comparer les aires des deux figures grisées du quadrillage.
\label{fig-ex-13-1}\figureGeo{aa-1}
\end{exercice}\end{blocLivre}
\begin{blocLivre}{ex-13-2}\begin{exercice}[1]{Mesurer en carreaux}
Donner l’aire de la figure en L, en carreaux.
\label{fig-ex-13-2}\figureGeo{aa-2}
\end{exercice}\end{blocLivre}
\begin{blocLivre}{ex-13-3}\begin{exercice}[1]{Formes différentes, aires égales}
Montrer que les deux figures ont la même aire.
\label{fig-ex-13-3}\figureGeo{aa-3}
\end{exercice}\end{blocLivre}
\begin{blocLivre}{ex-13-4}\begin{exercice}[1]{En centimètres carrés}
Chaque carreau du quadrillage est un carré de 1 cm de côté. Donner l’aire du rectangle.
\label{fig-ex-13-4}\figureGeo{aa-4}
\end{exercice}\end{blocLivre}
\begin{blocLivre}{ex-13-5}\begin{exercice}[1]{Nommer un angle}
Nommer l’angle marqué sur la figure, préciser son sommet et ses côtés.
\label{fig-ex-13-5}\figureGeo{aa-5}
\end{exercice}\end{blocLivre}

\rubrique{Entraînement}
\begin{blocLivre}{ex-13-6}\begin{exercice}[2]{Ranger des angles}
Ranger les trois angles dessinés du plus fermé au plus ouvert.
\label{fig-ex-13-6}\figureGeo{aa-6}
\end{exercice}\end{blocLivre}
\begin{blocLivre}{ex-13-7}\begin{exercice}[2]{Chercher les angles droits}
Avec une équerre, repérer les angles droits du quadrilatère dessiné.
\label{fig-ex-13-7}\figureGeo{aa-7}
\end{exercice}\end{blocLivre}
\begin{blocLivre}{ex-13-8}\begin{exercice}[2]{Des moitiés de carreaux}
Donner l’aire de la figure, qui contient des carreaux entiers et des demi-carreaux.
\label{fig-ex-13-8}\figureGeo{aa-8}
\end{exercice}\end{blocLivre}
\begin{blocLivre}{ex-13-9}\begin{exercice}[2]{Même tour, autre surface}
Vérifier que les deux rectangles ont le même périmètre, puis comparer leurs aires.
\label{fig-ex-13-9}\figureGeo{aa-9}
\end{exercice}\end{blocLivre}
\begin{blocLivre}{ex-13-10}\begin{exercice}[2]{L’angle du calque}
Un angle décalqué se superpose exactement sur un angle de la figure. Lequel ? Justifier.
\label{fig-ex-13-10}\figureGeo{aa-10}
\end{exercice}\end{blocLivre}

\rubrique{Problèmes}
\begin{blocLivre}{ex-13-11}\begin{exercice}[3]{Les deux parcelles}
Sur le plan, chaque carreau représente 1 $\mathrm m^2$. Comparer les aires des parcelles P et Q, puis donner l’aire de la plus grande.
\label{fig-ex-13-11}\figureGeo{aa-11}
\end{exercice}\end{blocLivre}
\begin{blocLivre}{ex-13-12}\begin{exercice}[3]{La frise de triangles}
La frise est pavée de triangles tous identiques, chacun valant un demi-carreau. Donner son aire en carreaux.
\label{fig-ex-13-12}\figureGeo{aa-12}
\end{exercice}\end{blocLivre}
\begin{blocLivre}{ex-13-13}\begin{exercice}[3]{Le côté qui double}
Un carré a 2 carreaux de côté. On double son côté.\begin{questions}
\item Donner l’aire du petit et du grand carré.
\item L’aire a-t-elle doublé ?
\end{questions}
\label{fig-ex-13-13}\figureGeo{aa-13}
\end{exercice}\end{blocLivre}
\begin{blocLivre}{ex-13-14}\begin{exercice}[3]{Les angles du triangle}
Comparer chaque angle du triangle dessiné à l’angle droit de l’équerre.
\label{fig-ex-13-14}\figureGeo{aa-14}
\end{exercice}\end{blocLivre}
\begin{blocLivre}{ex-13-15}\begin{exercice}[3]{Grand tour, petite aire}
Milan affirme : « Plus le périmètre d’une figure est grand, plus son aire est grande. » Comparer le serpentin et le carré dessinés, puis répondre.
\label{fig-ex-13-15}\figureGeo{aa-15}
\end{exercice}\end{blocLivre}

\clearpage
